%% kapa: exam=1 round=23 year=2012 period=2 form=A subject=accounting subject_ko=회계학 seq=1 number=1
%% source: Q-Net 감정평가사 기출문제 (cst00309) artlSeq=5052580 file=raw/1008/5052580/2012년 제23회 감정평가사 제1차시험 기출문제(2교시 A형).pdf
%% conversion: kordoc 4.24.51 -> tools/gen1.py (auto); review=proofread-sonnet; flags=none
\begin{question}{1}
재무보고를 위한 개념체계의 ‘유용한 재무정보의 질적 특성’에 관한 설명으로 옳지 않은 것은?
\begin{choices}
  \item 보강적 질적 특성은 정보가 목적적합하지 않거나 충실하게 표현되지 않으면, 개별적으로든 집단적으로든 그 정보를 유용하게 할 수 없다.
  \item 재무정보가 예측가치를 갖기 위해서 그 자체가 예측치 또는 예상치일 필요는 없다.
  \item 보강적 질적 특성을 적용하는 것은 어떤 규정된 순서를 따르지 않는 반복적인 과정으로, 때로는 하나의 보강적 질적 특성이 다른 질적 특성의 극대화를 위해 감소되어야 할 수도 있다.
  \item 중요성은 개별 기업 재무보고서 관점에서 해당 정보와 관련된 항목의 성격이나 규모 또는 이 둘 모두에 근거하여 해당 기업에 특유한 측면의 목적적합성을 의미한다.
  \item 보강적 질적 특성에는 비교가능성, 검증가능성, 적시성 및 충실한 표현이 있다.
\end{choices}
\end{question}
\answer{5}
